% Source: https://artofproblemsolving.com/community/c3h1493225_some_3d_geo
% Stable identity: aops-topic-1493225-q01
Let the edge be $a=2017$, let $O$ be the centre of the tetrahedron, and let $r_0$ be its inradius (the distance from $O$ to each face). The height of the tetrahedron is $a\sqrt{\dfrac{2}{3}}$ and $O$ divides each height in the ratio $3:1$, so
\[
r_0=\dfrac{1}{4}\cdot a\sqrt{\dfrac{2}{3}}=\dfrac{a\sqrt6}{12},\qquad VO=3r_0=\dfrac{a\sqrt6}{4}
\]
for every vertex $V$.

\textbf{Centre of one sphere.} A sphere tangent to the three faces through $V$ has its centre $C_V$ on the axis $VO$ of that corner (by symmetry), and its radius is proportional to the distance from $V$. Since the point $O$ at distance $3r_0$ from $V$ is at distance $r_0$ from these faces, a sphere of radius $\rho$ has $VC_V=3\rho$, so $OC_V=3r_0-3\rho$.

\textbf{Tangency of the spheres.} The four centres are the images of the vertices under the homothety with centre $O$ and ratio $\dfrac{OC_V}{OV}=\dfrac{r_0-\rho}{r_0}$, so they form a regular tetrahedron with edge $a\cdot\dfrac{r_0-\rho}{r_0}$. The spheres are tangent when this edge equals $2\rho$:
\[
a\left(1-\dfrac{\rho}{r_0}\right)=2\rho\quad\Longrightarrow\quad a=\rho\left(2+\dfrac{a}{r_0}\right)=\rho\left(2+2\sqrt6\right),
\]
since $\dfrac{a}{r_0}=\dfrac{12}{\sqrt6}=2\sqrt6$. Hence
\[
\rho=\dfrac{a}{2\left(1+\sqrt6\right)}=\dfrac{a\left(\sqrt6-1\right)}{10}.
\]
The figure shows the cross-section through the edge $VW$ and the midpoint $M$ of the opposite edge, a plane of symmetry that contains $O$ and the centres $C_V$, $C_W$ (the circles are great circles of two of the spheres, each tangent to the line of a face height).

\begin{asy}
usepackage("amsmath");
size(6.5cm);
defaultpen(fontsize(10pt));
real a=1; real hm=a/sqrt(2);
pair V=(-a/2,0),W=(a/2,0),M=(0,-hm),O=(0,-hm/2);
real R=a*sqrt(6)/4, rho=a*(sqrt(6)-1)/10;
pair CV=V+(3rho/R)*(O-V), CW=W+(3rho/R)*(O-W);
draw(V--W--M--cycle,linewidth(0.9));
draw(V--O--W,gray+dashed);
draw(circle(CV,rho),blue); draw(circle(CW,rho),blue);
draw(CV--CW,red); draw(CW--CW+rho*dir(70),deepgreen);
dot(V);dot(W);dot(M);dot(O);dot(CV,blue);dot(CW,blue);
label("$V$",V,NW);label("$W$",W,NE);label("$M$",M,S);label("$O$",O,S);
label("$\rho$",CW+0.5*rho*dir(70),dir(160),deepgreen+fontsize(8pt));
label("$2\rho$",CV+0.35*(CW-CV),N,red+fontsize(8pt));
label("\mbox{}",(min(currentpicture,user=true)+max(currentpicture,user=true))/2,black); // keeps the text after the figure black when LaTeX typesets the labels inline
\end{asy}

Each sphere also fits: its centre is at distance $4r_0-3\rho$ from the fourth face, which is more than $\rho$ because $\rho<r_0$. For $a=2017$ the radius is
\[
\boxed{\dfrac{2017\left(\sqrt6-1\right)}{10}}\approx292.4 .
\]
